\chapter{The Bladder}
\index{The Bladder}

\section*{Definition and Overview}
The bladder is a hollow, muscular organ situated in the lower pelvis that stores urine produced by the kidneys and releases it through the urethra. Urination is normally under voluntary control, coordinated by the bladder muscle (detrusor), the pelvic floor, the urethral sphincters, and nerves that link the bladder to the brain.

Bladder problems are common and include infection, overactivity, incontinence, stones, obstruction, and cancer. Symptoms such as burning on urination, frequency, blood in the urine, or difficulty emptying the bladder usually point to a disorder that can be diagnosed and effectively treated.

\section*{Epidemiology}
Bladder conditions are extremely common across all ages.

\begin{itemize}
    \item \textbf{Urinary tract infections:} affect women far more often than men, and roughly half of all women have at least one in their lifetime
    \item \textbf{Overactive bladder and incontinence:} become more common with age and affect both sexes, although stress incontinence is more common in women, often after childbirth
    \item \textbf{Bladder stones:} occur more often in men, particularly in association with an enlarged prostate
    \item \textbf{Bladder cancer:} about three to four times more common in men than in women, and strongly linked to smoking
    \item \textbf{Overall burden:} up to one in four older adults has some form of bladder problem
\end{itemize}

\section*{Aetiology and Pathophysiology}
Bladder problems arise from infection, muscle or nerve dysfunction, obstruction, or abnormal growth.

\begin{itemize}
    \item \textbf{Infection:} bacteria, usually from the bowel, travel up the urethra and multiply in the bladder, causing inflammation known as cystitis
    \item \textbf{Overactivity:} the detrusor muscle contracts too early or too strongly, producing urgency and frequency
    \item \textbf{Weakness:} weakened pelvic floor muscles or a weak bladder neck allow leakage with coughing, laughing, or exertion
    \item \textbf{Obstruction:} an enlarged prostate, a narrowed urethra, or severe constipation can block the flow of urine
    \item \textbf{Incomplete emptying:} a bladder that does not empty fully allows urine to stagnate, increasing the risk of infection and stone formation
    \item \textbf{Nerve damage:} diabetes, spinal injury, and some neurological conditions disrupt the signalling that controls the bladder
    \item \textbf{Stone formation:} concentrated urine and incomplete emptying allow crystals to aggregate into stones
    \item \textbf{Cancer:} smoking, exposure to certain chemicals, and chronic irritation can cause abnormal cell growth in the bladder lining
\end{itemize}

\section*{Clinical Features}
\begin{itemize}
    \item \textbf{Cystitis:} burning or stinging on urination, needing to pass urine frequently and urgently, and cloudy or strong-smelling urine
    \item \textbf{Blood in the urine:} visible or microscopic blood, which is a key warning sign and always requires investigation
    \item \textbf{Overactive bladder:} a sudden, strong urge to urinate, frequent trips to the toilet, and waking at night to urinate
    \item \textbf{Incontinence:} leaking urine with a cough, laugh, or exercise (stress incontinence), or a sudden urge with leakage (urge incontinence)
    \item \textbf{Obstruction:} a weak urine stream, hesitancy, straining, dribbling, and a sensation of incomplete emptying
    \item \textbf{Pain:} lower abdominal or pelvic pain, or discomfort when the bladder is full
    \item \textbf{Inability to urinate:} acute urinary retention, which is a medical emergency
    \item \textbf{Fever and back pain:} may indicate that the infection has spread to the kidneys
\end{itemize}

\section*{Differential Diagnosis}
\begin{itemize}
    \item Bladder infection versus kidney infection versus infection elsewhere in the urinary tract
    \item Overactive bladder versus stress incontinence versus mixed incontinence
    \item An enlarged prostate in men, which produces symptoms similar to bladder conditions
    \item Interstitial cystitis (bladder pain syndrome), which mimics recurrent infection
    \item Bladder stones versus other causes of blood in the urine
    \item Bladder cancer, which may present with painless blood in the urine
    \item Neurological conditions, such as spinal cord problems, causing bladder dysfunction
    \item Gynaecological causes in women, such as vaginitis or pelvic floor prolapse, and prostate problems in men
\end{itemize}

\section*{When to Seek Medical Care}
See a doctor for any blood in the urine, pain or burning with urination, urgency or frequency that persists, difficulty emptying the bladder, or leaking urine that bothers you. People with diabetes, a history of bladder problems, or a weakened immune system should be assessed early. Anyone unable to pass urine, or passing only small amounts with severe pain, needs urgent care.

\textbf{Contact your local emergency services for:}
\begin{itemize}
    \item Inability to pass urine with severe lower abdominal pain, indicating urinary retention
    \item Fever, chills, and back pain suggesting a kidney infection
    \item Visible blood in the urine with the formation of clots or difficulty passing urine
    \item Sudden loss of bladder control with back or leg weakness, which may indicate a spinal problem
\end{itemize}

\section*{Diagnosis and Investigations}
\begin{itemize}
    \item \textbf{Urine test:} a dipstick and laboratory analysis of the urine to detect infection, blood, sugar, and protein
    \item \textbf{Urine culture:} to identify the bacteria and guide the choice of antibiotic
    \item \textbf{History:} review of fluid intake, medicines, and a symptom diary of urination
    \item \textbf{Examination:} abdominal and pelvic examination, including a prostate examination in men
    \item \textbf{Bladder scan:} ultrasound to measure the volume of urine remaining after voiding
    \item \textbf{Urodynamic studies:} measurements of bladder pressure and flow to clarify how the bladder is functioning
    \item \textbf{Cystoscopy:} a small camera passed into the bladder to inspect its lining
    \item \textbf{Imaging:} ultrasound, CT, or MRI of the kidneys and bladder
    \item \textbf{Blood tests:} kidney function and, where relevant, prostate-specific antigen in men
\end{itemize}

\section*{Management}
\begin{itemize}
    \item \textbf{Infection:} a short course of antibiotics, with generous fluid intake to flush the bladder
    \item \textbf{Bladder retraining:} timed voiding and techniques to delay the urge to urinate
    \item \textbf{Pelvic floor exercises:} strengthening exercises that improve stress incontinence in many women and some men
    \item \textbf{Medicines:} anticholinergic and other agents that calm an overactive bladder
    \item \textbf{Obstruction:} treatment of the underlying cause, such as medicines or surgery for an enlarged prostate
    \item \textbf{Stones:} may pass naturally with fluids, or be removed by endoscopic procedures
    \item \textbf{Cancer:} surgical removal of the tumour, often combined with other treatments under specialist care
    \item \textbf{Procedures:} injections into the bladder wall, nerve stimulation, or surgery for severe incontinence
    \item \textbf{Catheters:} used in the short or long term when the bladder cannot empty on its own
    \item \textbf{Allied health:} physiotherapy and continence nursing support
\end{itemize}

\section*{Complications}
\begin{itemize}
    \item Recurrent infections and, if untreated, spread of infection to the kidneys (pyelonephritis)
    \item Kidney damage from long-standing obstruction or reflux of urine
    \item Bladder stones and chronic irritation of the bladder lining
    \item Worsening incontinence, skin breakdown, and reduced quality of life
    \item Social embarrassment, anxiety, and avoidance of activity
    \item Delays in diagnosing bladder cancer, which can reduce treatment options
    \item Falls in older people who rush to the toilet at night
\end{itemize}

\section*{Prognosis and Outcomes}
Simple bladder infections clear quickly with antibiotics in the great majority of cases. Overactive bladder and incontinence respond well to bladder retraining, pelvic floor exercises, and medicines in many people, although some require long-term treatment.

Bladder stones are usually cured once removed and the underlying cause is addressed. The outlook for bladder cancer depends on how early it is detected and its stage, and many early cancers are curable. Early investigation of blood in the urine offers the best chance of a good outcome.

\section*{Prevention and Self-Care}
\begin{itemize}
    \item Drink enough fluid through the day to keep the urine dilute, and urinate when you feel the need
    \item Empty the bladder fully, and for women urinate shortly after sexual activity to reduce the risk of infection
    \item Wipe front to back after using the toilet to reduce the spread of bacteria
    \item Avoid constipation by eating fibre, drinking fluids, and staying active
    \item Do pelvic floor exercises regularly, especially after childbirth
    \item Maintain a healthy weight and limit caffeine, alcohol, and carbonated drinks that irritate the bladder
    \item Do not smoke, because smoking greatly increases the risk of bladder cancer
    \item See a doctor promptly for any blood in the urine, even if it stops by itself
\end{itemize}

\section*{Sources and Further Reading}
The following organisations provide reliable information on bladder health and bladder disorders for patients, families, and health professionals.

\begin{itemize}
    \item NHS. \textit{Information on bladder and urinary tract conditions.} https://www.nhs.uk/
    \item Mayo Clinic. \textit{Overview of bladder disorders, including infections and incontinence.} https://www.mayoclinic.org/
    \item MedlinePlus. \textit{Health information on bladder diseases and urinary problems.} https://medlineplus.gov/
    \item NICE. \textit{Clinical guidance on the management of lower urinary tract symptoms and incontinence.} https://www.nice.org.uk/
    \item Centers for Disease Control and Prevention. \textit{Resources on urinary tract health and infection prevention.} https://www.cdc.gov/
\end{itemize}
